\subsection{Ensayo técnico: análisis de las preguntas i--n}

El análisis de la PSD en GNU Radio permite relacionar directamente la tasa de bits, la frecuencia de muestreo y las operaciones internas del flujograma con la forma observable del espectro. Si la señal binaria rectangular tiene una tasa de bits \(R_b\) y el sistema utiliza \(Sps\) muestras por símbolo, entonces la frecuencia de muestreo en el punto donde ya se ha realizado la interpolación es

\[
f_s = R_b \cdot Sps.
\]

Como el analizador de espectros representa el eje frecuencial de forma simétrica alrededor del cero, el rango total de frecuencias ocupado por la visualización de la PSD es

\[
\left[-\frac{f_s}{2},\frac{f_s}{2}\right]
=
\left[-\frac{R_b Sps}{2},\frac{R_b Sps}{2}\right].
\]

De esta manera se responde cómo calcular todo el rango de frecuencias cuando se conocen \(R_b\) y \(Sps\).

La resolución espectral del analizador depende del número de puntos de la FFT y de la frecuencia de muestreo. Si el bloque utiliza una FFT de \(N\) puntos, la separación entre líneas espectrales consecutivas es

\[
\Delta f = \frac{f_s}{N}.
\]

Esto significa que un mayor número de puntos en la FFT mejora la resolución en frecuencia, mientras que una frecuencia de muestreo más alta aumenta el rango visible del espectro para un mismo valor de \(N\).

En cuanto al bloque \texttt{Unpack K Bits}, su función es convertir cada muestra de entrada en \(K\) bits de salida. Si la entrada corresponde a bytes de un archivo, el valor natural es \(K=8\), puesto que cada byte contiene ocho bits. Si se configurara \(K=16\), la tasa de salida del bloque se duplicaría con respecto al caso normal, es decir,

\[
f_{s,\text{out}}^{(\text{Unpack})} = K\,f_{s,\text{in}}^{(\text{Unpack})},
\]

pero además se perdería coherencia con la estructura original del archivo, ya que se estaría forzando una descomposición de 16 bits sobre datos almacenados de manera natural en palabras de 8 bits. En la práctica, esto alteraría la interpretación correcta de la información y modificaría artificialmente la velocidad del flujo de datos.

Si se conoce el número de lóbulos observados en la PSD y el ancho de banda característico asociado a cada lóbulo, puede estimarse primero la frecuencia de muestreo visible en el analizador. Si se denota por \(L\) el número de lóbulos contados según la convención del laboratorio y por \(B\) el ancho asociado a cada lóbulo, entonces una aproximación útil es

\[
f_{s,\text{PSD}} \approx L \, B.
\]

Una vez obtenida esa frecuencia de muestreo en el punto de observación, es posible retroceder por el flujograma para hallar la frecuencia a la entrada del bloque \texttt{Unpack K Bits}. Si la señal observada ya pasó por el \texttt{Unpack K Bits} y luego por una interpolación de factor \(Sps\), entonces

\[
f_{s,\text{in}}^{(\text{Unpack})}
=
\frac{f_{s,\text{PSD}}}{K \cdot Sps}.
\]

A partir de esta expresión se deduce también la frecuencia de muestreo a la salida del propio bloque \texttt{Unpack K Bits}, que simplemente queda dada por

\[
f_{s,\text{out}}^{(\text{Unpack})}
=
K\,f_{s,\text{in}}^{(\text{Unpack})}.
\]

Por su parte, el bloque \texttt{Char to Float} no cambia la tasa de muestreo. Su función es únicamente convertir el tipo de dato para que las muestras puedan ser procesadas en formato numérico de punto flotante por los bloques siguientes. Por ello, si se conoce la frecuencia de muestreo a su entrada, la frecuencia a su salida permanece igual:

\[
f_{s,\text{out}}^{(\text{Char to Float})}
=
f_{s,\text{in}}^{(\text{Char to Float})}.
\]

En síntesis, el rango frecuencial visible en la PSD lo fija la frecuencia de muestreo, la resolución del analizador depende de la FFT, el bloque \texttt{Unpack K Bits} multiplica la tasa por \(K\), y el bloque \texttt{Char to Float} conserva dicha tasa. Estas relaciones permiten interpretar correctamente el comportamiento del espectro observado en GNU Radio y justifican matemáticamente los cambios producidos al modificar la fuente de datos o los parámetros del flujograma.